\chapter{Estado del arte}
\label{ch:estado_arte}

\section{Riesgo de inundación en un contexto cambiante}

Las inundaciones constituyen el desastre natural con mayor frecuencia, mayor extensión geográfica y mayor impacto económico a nivel global \parencite{najibi2018global}. En África y otras regiones en desarrollo, su impacto social es especialmente severo y con frecuencia subestimado \parencite{dibaldassarre2010}. En un contexto marcado por el cambio climático, el crecimiento urbano y la expansión de la exposición, los sistemas de análisis tradicionales, basados en escenarios deterministas y períodos de retorno de la variable forzante, resultan insuficientes para representar la variabilidad natural y la incertidumbre que caracterizan los eventos extremos.

Los enfoques probabilísticos y estocásticos ofrecen una representación más completa, pero su adopción práctica sigue siendo limitada. La curva de aprendizaje técnica, la falta de herramientas integradas y la complejidad operativa de las cadenas de cálculo impiden que muchos equipos profesionales los utilicen de forma sistemática \parencite{dibaldassarre2010}. La guía metodológica española para cartografía de zonas inundables reconoce esta tensión entre rigor metodológico y aplicabilidad práctica \parencite{marm2011guia}, lo que refuerza la pertinencia de una herramienta como \hydra.

\section{Modelado estocástico de inundaciones}
\label{sec:modelado_estocastico}

El modelado de inundaciones ha estado dominado históricamente por enfoques deterministas centrados en escenarios de diseño con un único evento representativo. Sin embargo, dos eventos con el mismo período de retorno en términos de precipitación o caudal pueden producir impactos hidráulicos muy diferentes en función de la distribución temporal de la lluvia, la humedad antecedente, la respuesta de la cuenca o la geometría hidráulica local. El período de retorno del forzamiento no es el período de retorno del daño.

Los enfoques estocásticos amplían el espacio de simulación generando múltiples eventos plausibles y calculando la distribución de impactos resultante. Esta distinción conceptual, del período de retorno de la precipitación o el caudal al período de retorno del impacto hidráulico, es uno de los pilares metodológicos de esta tesis \parencite{navas2017jia,navas2018besaya,navas2024calle30}. Para lograrlo de forma operativa, la metodología combina generación sintética de eventos, modelos hidrológicos e hidráulicos, selección de escenarios representativos y técnicas de reconstrucción o emulación para cubrir el espacio de impactos.

En la práctica, el reto consiste en mantener un equilibrio entre realismo físico y coste computacional. Las metodologías modernas suelen estructurarse en tres etapas: generación masiva de forzamientos sintéticos, selección mediante algoritmos como la máxima disimilitud, y simulación completa del subconjunto seleccionado seguida de reconstrucción estadística del resto \parencite{ihcantabria2020foresee,ihcantabria2021m30}. Este esquema ha sido aplicado en contextos reales de inundación fluvial urbana y en estudios de resiliencia de infraestructuras críticas.

\section{Análisis estadístico de extremos e inferencia bayesiana}
\label{sec:extremos}

El análisis de extremos proporciona el marco estadístico para estimar cuantiles raros y períodos de retorno. Las formulaciones clásicas, basadas en distribuciones GEV y modelos POT con máxima verosimilitud, han sido extendidas con métodos bayesianos que incorporan incertidumbre paramétrica, conocimiento a priori e incertidumbre estructural \parencite{martins2000generalized,coles2001extremes}.

Los modelos bayesianos jerárquicos permiten aprovechar simultáneamente información de múltiples estaciones o períodos, mejorando la estimación en cuencas con registros cortos \parencite{cooley2007bayesian}. Herramientas como Stan \parencite{carpenter2017stan}, PyMC \parencite{salvatier2016pymc}, R-INLA \parencite{rue2009inla} y OpenTURNS \parencite{baudin2016openturns} proporcionan soporte computacional robusto para estos análisis. El paquete extRemes en R implementa metodologías clásicas de análisis de extremos en un entorno accesible \parencite{gilleland2016extremes}. Sin embargo, su adopción en entornos aplicados exige encapsular complejidad y estandarizar flujos, lo que constituye uno de los objetivos de \hydra.

El desplazamiento del foco desde los extremos de la variable forzante hacia los extremos del impacto hidráulico es una idea central de esta tesis. La incertidumbre en parámetros de distribuciones de extremos, especialmente en series cortas o eventos raros, puede ser tan relevante como la estimación central. El episodio de precipitación extrema de Valencia de octubre de 2024 ilustra la importancia de este enfoque: la estimación de niveles de retorno de un evento sin precedente histórico exige tanto un marco estadístico bayesiano robusto como series de referencia adecuadas \parencite{deljesus2025valencia}.

\section{Análisis de frecuencia regional y geoestadística}
\label{sec:rfa}

Cuando las series locales son cortas o los cuantiles de interés corresponden a períodos de retorno superiores a la longitud del registro, el análisis de frecuencia regional (RFA) permite combinar información de múltiples estaciones bajo la hipótesis de homogeneidad regional. El enfoque basado en L-momentos de \textcite{hosking1997rfa} es el más extendido: permite identificar regiones homogéneas, seleccionar distribuciones regionales y estimar cuantiles con mayor precisión que los métodos locales cuando el número de años de registro es limitado.

La interpolación espacial de parámetros de distribución y de cuantiles extremos es una extensión natural del RFA. Las técnicas de kriging y los modelos geoestadísticos permiten mapear incertidumbre y aprovechar estructuras de correlación espacial que los enfoques puramente locales ignoran \parencite{cooley2007bayesian}. En \hydra, el módulo de análisis espacial implementa estas capacidades, permitiendo tanto el análisis regional de frecuencias como la generación de campos espaciales de extremos consistentes con la estructura de dependencia observada.

\section{Análisis multivariante y cópulas}
\label{sec:cópulas}

Los extremos hidrológicos y sus impactos tienen carácter multivariante. La intensidad de precipitación, la duración del evento, la humedad antecedente del suelo, la forma del hidrograma y la altura de la llanura de inundación son variables que interactúan y cuya dependencia condicionada determina el impacto final. Tratar estas variables como independientes conduce a estimaciones sesgadas del riesgo.

Las cópulas proporcionan un marco matemático flexible para modelar estructuras de dependencia entre variables aleatorias sin imponer restricciones sobre sus distribuciones marginales \parencite{genest2007copulas,salvadori2004copulas,salvadori2007extremes}. En el contexto de la generación de eventos sintéticos para estudios de inundación, las cópulas gaussianas permiten generar conjuntos de eventos que preservan la correlación multivariante observada en los registros históricos, manteniendo al mismo tiempo la estructura de márgenes individuales ajustada mediante análisis de extremos. Esta combinación, cópulas para la dependencia y distribuciones GEV o POT para las márgenes, es la base del generador de eventos de caudal empleado en Besaya y del generador de eventos de precipitación multisitio empleado en Calle 30 \parencite{navas2018besaya,navas2024calle30}.

El análisis multivariante basado en cópulas constituye uno de los componentes metodológicos relevantes incorporados en \pyhydra frente a los enfoques tradicionales univariantes. La comparación entre ambas aproximaciones, realizada en el caso de Calle 30, demuestra que los enfoques univariantes tienden a subestimar la frecuencia de eventos extremos al ignorar la dependencia entre variables \parencite{navas2024calle30}.

Una limitación de las cópulas elípticas es su incapacidad para representar dependencias asimétricas o de cola pesada, frecuentes en precipitación extrema multisitio. Las vine cópulas \parencite{aas2009vines} superan esta restricción descomponiendo la densidad multivariante en un producto de cópulas bivariantes condicionadas organizadas en una estructura jerárquica de árboles. Este enfoque, aplicado por \textcite{urrea2026vine} al río Besaya, demuestra que las vine cópulas flexibles superan sistemáticamente a la cópula gaussiana en la reproducción de la dependencia conjunta en la cola superior de la distribución, y que el cálculo de períodos de retorno multivariantes puede acelerarse en un orden de magnitud mediante la aproximación de la función de distribución conjunta con procesos gaussianos (GPR).

\section{Generación estocástica de precipitación}
\label{sec:generacion_estocastica}

La generación sintética de series de precipitación tiene una larga historia teórica que arranca en los modelos de puntos de Poisson y sus extensiones con estructura de cluster. Los procesos de Neyman-Scott de pulsos rectangulares (NSRP), formulados originalmente por \textcite{rodriguezituribe1987,rodriguezituribe1988}, describen la precipitación como una superposición de tormentas, cada una asociada a un evento de Poisson con un número aleatorio de celdas que producen pulsos de intensidad y duración constantes. Estos modelos, calibrables a partir de estadísticos de primer y segundo orden de series observadas, reproducen propiedades esenciales como la distribución de intensidades, la proporcion de períodos secos, la autocorrelacion temporal y la distribución de acumulados a distintas escalas. La extensión espacial STNSRP añade dependencia entre estaciones y permite generar campos multisitio coherentes.

NEOPRENE no es un proceso estocástico adicional, sino la librería Python que contiene la implementación NSRP puntual y STNSRP multisitio, con soporte para generación espacial de lluvia y desagregación temporal de acumulados diarios a resoluciones subhorarias \parencite{diezsierra2023neoprene}. Su diseño como paquete independiente, documentado y distribuido en GitHub y Zenodo, establece un precedente directo para la arquitectura modular de \pyhydra. La integración de NEOPRENE como backend de \path{NSRPModel} y \path{STNSRPModel} en el submódulo \path{climate.stochastic_generation} es una de las conexiones más directas entre publicaciones previas y los módulos del bloque de análisis climático-estadístico.

Además de la ruta NEOPRENE (NSRP/STNSRP), \pyhydra incorpora dos rutas complementarias. La primera es un generador puntual de eventos de precipitación basado en separación de eventos y modelado con cópulas gaussianas, útil para construir escenarios discretos de simulación hidrológica directa \parencite{navas2018besaya,navas2024calle30}. La segunda es la integración de CoSMoS como librería de generación de series hidroclimáticas generales, orientada a preservar distribución marginal, dependencia temporal y estacionalidad. La distinción es importante: NSRP/STNSRP, implementados en NEOPRENE, representan explícitamente tormentas y células de lluvia; CoSMoS proporciona un marco estocástico general para series temporales; y el generador por cópulas se orienta a eventos discretos.

\section{Cambio climático: datos, proyecciones y corrección de sesgo}
\label{sec:cambio_climatico}

La evaluación del riesgo de inundación bajo escenarios futuros exige integrar proyecciones climáticas de modelos de circulación general o regional. El programa CMIP6 \parencite{eyring2016cmip6}, junto con los escenarios SSP \parencite{oneill2016ssp}, proporciona el marco de referencia actual para explorar la incertidumbre intermodelo y multiescenario. El acceso a estos datos se realiza a través de plataformas como la ESGF y el Copernicus Climate Data Store, que \hydra utiliza mediante sus módulos de fuentes de datos.

Los modelos climáticos presentan sesgos sistemáticos respecto a las observaciones locales que deben corregirse antes de emplear sus salidas en estudios hidrológicos. Los métodos de corrección de sesgo incluyen el escalado delta, el mapeo cuantil y sus variantes que preservan la señal de cambio climático \parencite{gudmundsson2012downscaling,teutschbein2012biascorrection}. La combinación de proyecciones de múltiples modelos y escenarios permite caracterizar la incertidumbre intermodelo, que en estudios de inundación puede ser comparable o superior a la incertidumbre estadística de los extremos \parencite{cloke2013modelling,smith2014comparing}.

El trabajo sobre los efectos del cambio climático en los recursos hídricos de los países andinos \parencite{deljesus2020andinos}, realizado con el modelo VIC y proyecciones CMIP5, y el posterior informe para el Banco Interamericano de Desarrollo \parencite{paz2019idb} son los antecedentes directos del bloque de cambio climático de \pyhydra. La herramienta SIMPCCe, que implementa flujos automatizados de descarga, corrección de sesgo y análisis de aportaciones a embalses para cuencas españolas, es el precedente más reciente de este enfoque transferido a una herramienta Python utilizable \parencite{navas2025simpce}.

\section{Fuentes de datos: reanálisis, teledetección y proyecciones}
\label{sec:fuentes_datos_arte}

El reanálisis ERA5 de ECMWF \parencite{hersbach2020era5} proporciona series horarias de precipitación, temperatura, viento y otras variables atmosféricas con cobertura global y resolución espacial de 0.25 grados desde 1940. Es la referencia principal para estudios climáticos e hidrológicos cuando no existen observaciones locales suficientes. Su acceso a través del Copernicus Climate Data Store permite descargas programaticas que \hydra automatiza en su módulo COPERNICUS.

Para precipitación de alta resolución espacial, los productos satelitales GPM-IMERG \parencite{huffman2020gpm} y PERSIANN-CCS \parencite{hong2004persiannccs} ofrecen estimaciones subdiarias con cobertura semiglobal. Otros productos de la misma familia, como PERSIANN-CDR \parencite{ashouri2015persiann}, aportan registros diarios de mayor longitud temporal para análisis climatológicos. Estos productos son especialmente útiles en regiones sin red densa de pluviómetros y en el análisis de eventos convectivos de alta intensidad. Su sesgo sistemático respecto a observaciones in situ exige aplicar técnicas de corrección antes de emplearlos en análisis de extremos o calibración hidrológica.

Los caudales observados a escala global se encuentran disponibles en productos como GloFAS \parencite{alfieri2013glofas} y la base GRDC. GloFAS combina el modelo hidrológico LISFLOOD con el sistema de predicción de conjunto de ECMWF para producir series de caudal en puntos de la red fluvial global, lo que lo convierte en una fuente especialmente útil en cuencas sin aforos o para validación de largo plazo.

Para la parametrización de modelos hidrológicos, la información edafológica de SoilGrids \parencite{poggio2021soilgrids} proporciona texturas, clases texturales y propiedades físico-hidrológicas a 250 metros de resolución con estimación de incertidumbre, lo que permite una parametrización más robusta que los productos de suelos tradicionales en regiones con coberturas edafológicas incompletas.

\section{Automatización de modelos hidrológicos e hidráulicos}
\label{sec:automatización}

Los modelos HEC-HMS \parencite{scharffenberg2016hechms}, SWAT \parencite{arnold1998swat,gassman2007swat}, SFINCS \parencite{leijnse2021sfincs} y HEC-RAS \parencite{brunner2016hecras} constituyen el núcleo de la simulación física en \hydra. Cada uno cuenta con una larga trayectoria técnica y una comunidad de usuarios consolidada, pero su uso combinado en simulaciones estocásticas masivas sigue requiriendo una carga operativa elevada.

HEC-HMS y HEC-RAS disponen de interfaces gráficas potentes pero sus versiones automatizables requieren el uso de scripting en HEC-DSS, RAS Controller o el API JSON de HEC-RAS, que no están integrados entre sí. SWAT dispone de herramientas como SWAT+ Editor y SWAT-CUP, pero sigue dependiendo de archivos de texto y procesos secuenciales difíciles de escalar. SFINCS, aunque eficiente para simulaciones masivas bidimensionales, requiere conocimiento de estructuras de malla, formatos NetCDF y control de parámetros a través de scripts específicos.

La experiencia con HydroMT \parencite{eilander2023hydromt} ha demostrado el valor de construir modelos reproducibles mediante interfaces programaticas. Inspirado en este precedente, \hydra adopta una arquitectura de adaptadores: cada modelo se encapsula en un módulo que traduce datos internos de la librería a los formatos esperados, lanza la ejecución y recupera y homogeniza las salidas. Esto permite organizar flujos de simulación masiva sin alterar los motores físicos ni su aceptacion técnica por parte de administraciones y consultoría.

Los casos de aplicación del proyecto FORESEE \parencite{ihcantabria2020foresee,ihcantabria2021m30} y los estudios realizados para el Besaya, Mallorca y Calle 30 \parencite{navas2018besaya,navas2019mallorca,navas2024calle30} han puesto de manifiesto los puntos exactos de friccion en la automatización de cada modelo y han informado las decisiones de diseño de los adaptadores implementados en \pyhydra.

\section{Downscaling híbrido y reconstrucción de manchas de inundación}
\label{sec:downscaling}

El downscaling híbrido combina información de modelos dinámicos con métodos estadísticos y de aprendizaje automático para reconstruir variables de alta resolución a partir de simulaciones de escala gruesa. En el contexto de inundaciones, esta estrategia permite estimar manchas de inundación para eventos no simulados directamente, aprovechando relaciones estadísticas aprendidas a partir de un subconjunto representativo de simulaciones completas.

La estrategia implementada en \hydra se basa en dos componentes. El primero es la selección de casos a simular mediante el algoritmo de máxima disimilitud, que garantiza que el subconjunto seleccionado cubre de forma uniforme el espacio de variabilidad de los eventos generados. El segundo es la reconstrucción de manchas de inundación para el resto de eventos mediante técnicas de k-vecinos más próximos (k-NN) o modelos estadísticos que interpolan en el espacio de características del evento \parencite{ihcantabria2020foresee,ihcantabria2021m30,navas2024calle30}. Esta combinación, denominada downscaling híbrido en la literatura más reciente \parencite{deljesus2024downscaling}, permite reducir el coste computacional del análisis estocástico sin sacrificar la cobertura del espacio de posibles impactos.

El enfoque conecta con desarrollos recientes que combinan generadores estocásticos de precipitación basados en el proceso de Neyman-Scott con técnicas de aprendizaje automático para mejorar la resolución espacial y temporal de los forzamientos hidrológicos \parencite{deljesus2024downscaling}. La integración de estos desarrollos en un flujo automatizado reproducible es uno de los objetivos específicos de la Tarea 2 del plan de investigación.

\section{Brecha identificada}
\label{sec:brecha}

La revisión del estado del conocimiento pone de manifiesto que el modelado estocástico de inundaciones, el análisis de extremos con incertidumbre y la integración de escenarios de cambio climático han avanzado de forma significativa en el plano metodológico. Herramientas como Stan, PyMC o OpenTURNS permiten realizar análisis probabilísticos complejos; modelos como HEC-HMS, SWAT, SFINCS o HEC-RAS ofrecen simulaciones detalladas del ciclo hidrológico y la dinámica de inundaciones; y productos globales como ERA5, GPM o CMIP6 proporcionan datos de entrada con cobertura universal.

Sin embargo, persisten importantes limitaciones que impiden la aplicación generalizada de estas metodologías en entornos operativos. La brecha principal no es la ausencia de métodos aislados, sino la falta de integración operativa entre datos climáticos, análisis estadístico, generación estocástica, modelos físicos y documentación utilizable. Cada uno de estos componentes requiere conocimientos técnicos específicos y consume tiempo en adaptación de formatos, conversión de unidades, organización de ficheros y comprobación de resultados.

En estudios estocásticos de inundación, donde el número de simulaciones puede ser de cientos o miles, la gestión manual de este flujo es inviable. La automatización no solo reduce errores, sino que mejora la trazabilidad, facilita la auditoría de resultados y permite transferir la metodología a nuevos equipos o territorios sin reconstruir toda la cadena de cálculo. Esta es la brecha que \hydra busca cerrar.

\section{Línea de investigación propia}
\label{sec:linea_propia}

Una parte sustancial del estado del arte empleado en esta tesis procede de una línea de trabajo desarrollada de forma ininterrumpida desde 2017 en IHCantabria. Esta trayectoria combina investigación aplicada, transferencia a proyectos reales y publicación de herramientas o metodologías reproducibles. En conjunto, las publicaciones de esta línea no solo justifican las decisiones de diseño de \hydra, sino que constituyen los antecedentes directos de sus módulos y flujos de trabajo.

\paragraph{Primera presentación de la metodología (2017).}
La presentación inicial de la metodología estocástica para estudios de inundación tuvo lugar en las V Jornadas de Ingeniería del Agua \parencite{navas2017jia}. En ese trabajo se describió la cadena metodológica basada en generación sintética de precipitación mediante técnicas geoestadísticas, simulación hidrológica e hidráulica y análisis estadístico de los impactos resultantes mediante minería de datos y regresión probabilística. Ese esquema conceptual prefigura directamente la arquitectura de \hydra.

\paragraph{Primera aplicación metodológica completa: río Besaya, Los Corrales de Buelna (2018).}
El artículo en la Revista de Obras Públicas representa la primera aplicación metodológica completa de la cadena estocástica a un caso real, tomando como objeto el río Besaya a su paso por Los Corrales de Buelna, Cantabria \parencite{navas2018besaya}. El trabajo fue distinguido con el Premio Nacional al mejor Trabajo de Fin de Master del Colegio de Ingenieros de Caminos. La contribución conceptual clave es el desplazamiento del foco desde el período de retorno del caudal hacia la distribución de los impactos hidráulicos como calados y extensiones.

En esta primera formulación, el flujo parte de registros de caudal observados: se separan eventos, se caracterizan mediante variables como pico, volumen y duración, se ajusta una estructura de dependencia multivariante y se generan hidrogramas sintéticos plausibles. Solo un subconjunto representativo se simula hidráulicamente, en aquel momento con Iber y con un modelo digital del terreno derivado de datos LiDAR, para reconstruir posteriormente la distribución espacial de impactos. La automatización de HEC-RAS se incorporó en trabajos posteriores; por tanto, el caso de Corrales no debe interpretarse como una aplicación inicial de NEOPRENE ni de HEC-RAS, sino como el origen del enfoque estocástico de impacto que después se generaliza en \hydra. El esquema metodológico se recupera y formaliza en la \cref{fig:besaya_esquema_caudal}.

\paragraph{Extensión a Sant Llorenç des Cardassar (2019).}
La aplicación en Mallorca evaluó la metodología en un contexto mediterráneo con cuencas pequeñas, alta intensidad convectiva y respuesta hidrológica rápida \parencite{navas2019mallorca}. Fue el primer trabajo de la línea en el que se aplicó la metodología de downscaling híbrido partiendo de lluvia: generación estocástica de eventos de precipitación, selección de escenarios representativos, simulación física del subconjunto seleccionado y reconstrucción estadística de las manchas no simuladas. El trabajo incorporó pluviómetros horarios del SIAR y productos satelitales de precipitación, en particular PERSIANN \parencite{ashouri2015persiann}, para analizar la coherencia espacial de la lluvia en una zona con información observacional limitada. También se evaluó TRMM/TMPA \parencite{huffman2007trmm} como fuente alternativa, pero la reconstrucción final se apoyó en kriging de los registros pluviométricos porque proporcionaba mayor control local sobre la estructura espacial del evento que los productos satelitales o de reanálisis. Esta extensión validó la portabilidad del enfoque y expuso necesidades de adaptación en la generación de eventos, la reconstrucción espacial de precipitación y la parametrización hidrológica que han influido en el diseño modular de \hydra.

\paragraph{Impactos del cambio climático en sistemas hidroeléctricos andinos (2019--2020).}
La participacion en el informe del Banco Interamericano de Desarrollo \parencite{paz2019idb} y el artículo derivado \parencite{deljesus2020andinos} aportaron una perspectiva complementaria sobre el cambio climático, con el modelo hidrológico distribuido VIC, proyecciones CMIP5 bajo RCP4.5/8.5 y reconstrucción espacial por kriging. El estudio requirió la calibración automática del modelo hidrológico sobre múltiples cuencas andinas, lo que llevo al uso de \texttt{spotpy} \parencite{houska2015spotpy} como librería de optimización (algoritmos PSO, DREAM y SCE-UA). La experiencia acumulada en ese proceso informo directamente el desarrollo del módulo de calibración automática incorporado en la clase \texttt{HMSModel} de \hydra, que expone la misma interfaz \texttt{spotpy} para calibración de modelos HEC-HMS. Las lecciones sobre manejo de escenarios múltiples, incertidumbre intermodelo y transferencia de resultados a usuarios técnicos informan asimismo el bloque de cambio climático de \hydra.

\paragraph{NEOPRENE: generación espacial de lluvia en Python (2021--2023).}
La necesidad de un generador estocástico de precipitación espacial, mantenible y distribuido como paquete Python, condujo al desarrollo de NEOPRENE \parencite{diezsierra2023neoprene}. NEOPRENE contiene tanto la implementación NSRP puntual como la extensión STNSRP multisitio, con soporte para desagregación temporal y diagnósticos estadísticos. Su publicación en Geoscientific Model Development establece un precedente de software científico reproducible que \hydra adopta como modelo de referencia. La integración de NEOPRENE como librería de generación estocástica dentro del bloque de análisis climático-estadístico de \pyhydra es una de las conexiones más directas entre publicaciones previas y la arquitectura de la librería.

\paragraph{Calle 30, Madrid: validación en una infraestructura crítica (2024).}
La aplicación al sistema de túneles M-30 en Madrid \parencite{navas2024calle30} es el caso de mayor exigencia técnica e industrial de esta línea. A diferencia del Besaya, el flujo metodológico parte de precipitación multisitio, no de caudales observados. Los eventos de lluvia se separan por pluviómetro, se caracterizan estadísticamente, se generan eventos sintéticos mediante cópulas gaussianas y se reconstruye su distribución espacial antes de transformarlos en caudales mediante modelado hidrológico. La cadena se completa con selección de escenarios representativos, simulación hidráulica con HEC-RAS 1D en el sistema de túneles y reconstrucción de los resultados no simulados mediante k-NN.

Este caso consolida el salto desde una metodología estocástica fluvial basada en caudal hacia una metodología urbana basada en precipitación, hidrología e hidráulica acopladas. La comparación con enfoques clásicos univariantes demuestra que estos subestiman la frecuencia de eventos extremos al ignorar la dependencia entre variables de la lluvia y de la respuesta hidrológica. Desarrollado en colaboración con Ferrovial y en el marco del proyecto FORESEE \parencite{ihcantabria2020foresee,ihcantabria2021m30}, Calle 30 es la referencia industrial más directa que justifica la transferencia de \hydra a entornos profesionales. Su esquema metodológico se presenta en la \cref{fig:m30_esquema_precipitacion}.

La comparación entre ambos trabajos define dos entradas metodológicas complementarias que atraviesan toda la tesis. En cuencas con información hidrométrica suficiente, la cadena puede arrancar en la generación estocástica de caudales y estimar directamente la frecuencia de impactos hidráulicos. En sistemas urbanos o infraestructuras donde la lluvia espacialmente distribuida controla la respuesta, la cadena debe arrancar en la precipitación, incorporar reconstrucción espacial, modelado hidrológico y selección de escenarios antes de la simulación hidráulica. \hydra se diseña precisamente para alojar ambas rutas bajo una arquitectura común: datos, eventos, generación estocástica, modelos físicos, selección representativa y reconstrucción de impactos.

\paragraph{SIMPCCe: aportaciones a embalses bajo cambio climático (2025).}
SIMPCCe es una herramienta Python para el análisis de aportaciones a embalses bajo escenarios de cambio climático, integrada con fuentes de datos oficiales españolas \parencite{navas2025simpce}. Su desarrollo es una nueva iteración del ciclo publicación-herramienta que esta tesis generaliza, y sirve como precedente de integración de fuentes de datos meteorológicos e hidrológicos nacionales en flujos automatizados. La dimensión de transferencia del trabajo fue reconocida en el ámbito de la Fundación Botín, lo que refuerza su valor como antecedente industrial; al no existir una publicación científica específica para ese reconocimiento, la evidencia principal se mantiene en el artículo técnico de SIMPCCe.

\paragraph{Valencia 2024: comparación de métodos de ajuste de extremos (2025).}
La comunicación presentada en las VIII Jornadas de Ingeniería del Agua \parencite{deljesus2025jia} compara métodos de ajuste de la distribución GEV (MLE, L-momentos e inferencia bayesiana) y análisis de frecuencia regional a distintas escalas espaciales, tomando como caso de estudio el evento de precipitación extrema del 29 de octubre de 2024 en la Comunidad Valenciana; un artículo con los resultados completos ha sido enviado a \textit{Ingeniería del Agua} \parencite{navas2025valenciaIA}. Los resultados cuantifican la sensibilidad de cada método a la inclusión o exclusión del evento extremo y concluyen que la inferencia bayesiana proporciona estimaciones más estables y con incertidumbre explícitamente comunicable. Paralelamente, el resumen publicado en la EGU General Assembly 2025 \parencite{deljesus2025valencia} presenta los niveles de retorno en el marco internacional.

\paragraph{Lago Tanganica: niveles de diseño bajo cambio climático (2025).}
La comunicación sobre el Lago Tanganica \parencite{navas2025tanganika} desarrolla una metodología integrada para estimar niveles de diseño de infraestructuras portuarias en ausencia de datos instrumentales continuos. La cadena combina reanálisis ERA5, reconstrucción hidrológica de la cuenca del Congo, altimetría satelital Hydroweb, regresión multivariante con AdaBoost para estimar caudales, K-means para capturar la no linealidad de la respuesta del lago, y downscaling estadístico de 19 modelos CMIP6 bajo SSP2-4.5 y SSP5-8.5. Los incrementos proyectados de nivel para distintos horizontes temporales y períodos de retorno proporcionan una base cuantitativa para el rediseño del puerto de Kalundu (Burundi). Este caso extiende la arquitectura de \hydra a sistemas lacustres en contextos de escasa instrumentación. Un artículo con los resultados completos ha sido enviado a \textit{Ingeniería del Agua} \parencite{navas2025tanganicaIA}.

\paragraph{Besaya: incertidumbre de rugosidad y efectos de estructura de modelo (2025).}
El artículo enviado a \textit{Environmental Modelling \& Software} \parencite{navas2025roughness} compara SFINCS y HEC-RAS mediante un ensemble emparejado de 995 realizaciones de rugosidad de Manning, generado por Monte Carlo sobre nueve clases de uso del suelo en la llanura de inundación del río Besaya para un evento de $T = 500$ años. Los valores medios de Manning explican menos del 2\% de la varianza de las salidas en ambos modelos, pero los dos motores producen estructuras de incertidumbre marcadamente distintas: HEC-RAS genera una distribución bimodal del área inundada, ausente en SFINCS, atribuible al rebase de un umbral topográfico de 60,1~m que abre una llanura secundaria de 7,4~ha. La persistencia de la bimodalidad bajo métricas alternativas y verificaciones de robustez indica que la configuración completa del modelo ---no solo las ecuaciones de gobierno--- determina la forma de la distribución de riesgo hidráulico. Este trabajo motiva el módulo \texttt{modeling.hydraulic.sensitivity} de \hydra y establece que la automatización masiva no es solo una mejora operativa, sino una condición necesaria para detectar efectos de estructura de modelo no visibles en simulaciones individuales.


\paragraph{Síntesis.}
La \cref{tab:publicaciones_modulos} resume la correspondencia entre cada publicación de la línea propia y los módulos de \hydra a los que contribuye metodológica o instrumentalmente.

\begin{longtable}{p{0.36\textwidth}p{0.07\textwidth}p{0.21\textwidth}p{0.24\textwidth}}
  \caption{Correspondencia entre publicaciones de la línea de investigación propia y módulos de \hydra. Los artículos propios de la tesis se indican en la primera sección.}
  \label{tab:publicaciones_modulos} \\
  \toprule
  Título & Año & Publicación & Módulos de \hydra \\
  \midrule
  \endfirsthead
  \toprule
  Título & Año & Publicación & Módulos de \hydra \\
  \midrule
  \endhead
  \multicolumn{4}{l}{\textit{Artículos propios de la tesis doctoral}} \\
  \midrule
  \textit{Application of a new methodology to improve the estimation of flood frequencies in Calle 30, Madrid} \parencite{navas2024calle30} & 2024 & \textit{Ingeniería del Agua} (publicado) & Flujo completo, HEC-HMS, HEC-RAS 1D, cópulas, k-NN \\[4pt]
  \textit{SIMPCCe: A tool for the analysis of reservoir inflows under climate change scenarios} \parencite{navas2025simpce} & 2025 & \textit{Ingeniería del Agua} (publicado) & Climate Change, Bias Correction, SWAT \\[4pt]
  \textit{Comparación de métodos de ajuste para la distribución de precipitaciones extremas: Análisis del evento de octubre 2024 en Valencia} \parencite{navas2025valenciaIA} & 2025 & \textit{Ingeniería del Agua} (enviado) & Extremos (GEV, MLE, MCMC), RFA, incertidumbre \\[4pt]
  \textit{Metodología innovadora para modelar el impacto del cambio climático en el Lago Tanganica con datos globales y regionales} \parencite{navas2025tanganicaIA} & 2025 & \textit{Ingeniería del Agua} (enviado) & ERA5, CMIP6, Bias Correction, ML, extremos lacustres \\[4pt]
  \textit{From roughness uncertainty to model-structure effects in 2D flood inundation models: a Monte Carlo comparison of SFINCS and HEC-RAS} \parencite{navas2025roughness} & 2025 & \textit{Environ. Model. Softw.} (enviado) & SFINCS, HEC-RAS, Monte Carlo, incertidumbre hidráulica, Besaya \\[4pt]
  \midrule
  \multicolumn{4}{l}{\textit{Otras publicaciones de la línea de investigación}} \\
  \midrule
  \textit{Análisis del riesgo de inundación mediante técnicas estadísticas avanzadas} \parencite{navas2017jia} & 2017 & Jornadas IA & Arquitectura general, flujo completo \\[4pt]
  \textit{Evaluación y análisis del riesgo de inundación del Río Besaya a su paso por Los Corrales de Buelna, Cantabria} \parencite{navas2018besaya} & 2018 & Jornadas IA & Eventos de caudal, cópulas, impactos hidráulicos \\[4pt]
  \textit{Evaluación de una metodología para estudios de inundación basada en técnicas estadísticas avanzadas. Aplicación en Sant Llorenç des Cardassar (Mallorca)} \parencite{navas2019mallorca} & 2019 & Jornadas IA & SIAR, PERSIANN/TRMM, kriging, cópulas, MaxDiss, k-NN \\[4pt]
  \textit{Vulnerabilidad al cambio climático y medidas de adaptación de los sistemas hidroeléctricos en los países andinos} \parencite{paz2019idb} & 2019 & Informe BID & Climate Change, Bias Correction, SWAT \\[4pt]
  \textit{Efectos del cambio climático en el recurso hídrico de los países andinos} \parencite{deljesus2020andinos} & 2020 & \textit{Ingeniería del Agua} & Climate Change, Bias Correction, SWAT \\[4pt]
  \textit{NEOPRENE v1.0.1: A Python library for generating spatial rainfall based on the Neyman-Scott process} \parencite{diezsierra2023neoprene} & 2023 & \textit{Geosci. Model Dev.} & Stochastic Generation, NSRP y STNSRP \\[4pt]
  \textit{Rainfall downscaling using stochastic generators and machine learning} \parencite{deljesus2024downscaling} & 2024 & EGU 2024 & Hybrid Downscaling \\
  \bottomrule
\end{longtable}

Por tanto, esta memoria no revisa el estado del arte desde una posición externa: organiza, formaliza e integra una trayectoria previa de investigación aplicada. \hydra se plantea como la consolidación industrial de esa experiencia en una herramienta modular para estudios de inundación bajo incertidumbre.

\section{Posicionamiento de la tesis frente al estado del arte}
\label{sec:posicionamiento}

La revisión presentada en este capítulo describe un campo metodológicamente maduro. Las distribuciones GEV y GPD, los estimadores L-momentos y bayesianos, los procesos NSRP/STNSRP, las librerías NEOPRENE y CoSMoS, las cópulas multivariantes y las vine cópulas, las técnicas de corrección de sesgo, el downscaling híbrido y los motores de simulación hidrológica e hidráulica son metodologías o herramientas con fundamento teórico sólido y validación extensa en la literatura. Las plataformas de datos ERA5, GPM, CMIP6 y SoilGrids proporcionan acceso sin precedentes a variables atmosféricas e hidrológicas a escala global.

Esta tesis no propone nuevos métodos en ninguno de estos ámbitos individuales. Su posicionamiento es distinto: demuestra que es posible integrar y operacionalizar este conjunto de metodologías dentro de un marco reproducible común. \hydra actúa como capa de orquestación que conecta fuentes de datos, análisis estadístico, generación estocástica y modelos físicos en una cadena auditable y transferible. La contribución original no reside en el desarrollo de nuevas técnicas aisladas, sino en la demostración de que la principal barrera para la adopción operativa de estas metodologías es superable mediante una arquitectura modular bien diseñada.

Este posicionamiento es coherente con la naturaleza industrial de la memoria y con la brecha identificada en la \cref{sec:brecha}: la fragmentación entre componentes impide la adopción generalizada de estas metodologías en entornos profesionales, y \hydra cierra esa brecha sin sustituir ni competir con los desarrollos existentes que integra.
